\documentclass[../../main.tex]{subfiles}
\begin{document}
\section{浸入子流形}\label{section:ch05-浸入子流形}

在\cref{example:ch04-figure-eight} 与\cref{example:ch04-dense-curve} 中,我们看到单射浸入的像未必是嵌入子流形。对这些像仍说些有意义的话是可能的;它们的局部性质与嵌入子流形相同,但它们的整体拓扑可能不同。本节我们引进比嵌入子流形更一般的一类子流形,它们正是单射浸入的像。

设 $M$ 是带边或不带边的光滑流形。$M$ 的一个\textbf{浸入子流形}是 $M$ 的子集 $S\subseteq M$,连同一个拓扑与一个光滑结构,使包含映射 $S\hookrightarrow M$ 是单射光滑浸入。

正如我们用拓扑嵌入与光滑嵌入的差别提醒自己的那样,这里也有一个重要的术语问题。我们并没有要求 $S$ 上的拓扑为子空间拓扑。(如果它确实是子空间拓扑,则 $S$ 就是嵌入子流形。)因此,浸入子流形远不只是一个可被视为流形的子集:它是一个光滑流形,并被以单射浸入的方式包含在 $M$ 中。

注意,每个嵌入子流形自动也是浸入子流形;故浸入子流形类包含嵌入子流形类作为特例。当我们要强调一个浸入子流形不一定是嵌入子流形时,有时称它为\textbf{非嵌入的}。

浸入子流形的余维数定义方式与嵌入子流形相同: $\operatorname{codim}S=\dim M-\dim S$。

下面的命题可由\cref{proposition:ch05-images-of-embeddings} 作相应修改推出。

\begin{proposition}[Images of Injective Immersions]\label{proposition:ch05-injective-immersion-images}
设 $M$ 是带边或不带边的光滑流形, $N$ 是光滑流形, $F:N\to M$ 是单射光滑浸入。令 $S=F(N)$,并赋予 $S$ 由 $F$ 诱导的微分同胚拓扑与光滑结构。则 $S$ 是 $M$ 的浸入子流形,且 $F$ 是从 $N$ 到 $S$ 的微分同胚。
\end{proposition}

\begin{proof}
对\cref{proposition:ch05-images-of-embeddings} 的证明作如下修改即可:现在 $F$ 是到其像上的同胚这一假设被替换为 $F$ 是单射浸入,于是\eqref{align:ch05-inclusion-factor} 中第二个映射只是浸入而非嵌入。
\end{proof}

由于 $S$ 的拓扑与光滑结构是由浸入子流形定义的一部分,谈论某集合 $S\subseteq M$ 是浸入子流形,就预先假定了我们至少在心里已经选定了一个拓扑与光滑结构。正如上面的命题表明的,最常见的情形是当我们从单射浸入 $F:N\to M$ 出发,并把 $S$ 定义为其像、再配上由 $F$ 诱导到 $S$ 上的拓扑与光滑结构时。

\begin{example}[Some Immersed Submanifolds]\label{example:ch05-immersed-examples}
\cref{example:ch04-figure-eight} 的八字曲线 $E\subseteq\mathbb{R}^2$ 是 $1$ 维浸入子流形(像集加上 $F$ 诱导的拓扑与光滑结构),但不是嵌入子流形。

\cref{example:ch04-dense-curve} 的稠密曲线 $\gamma(\mathbb{R})$ 是 $\mathbb{T}^2$ 的 $1$ 维浸入子流形,但不是嵌入子流形,因为其像在子空间拓扑下稠密,故它不能是 $\mathbb{T}^2$ 的嵌入子流形(它是 $1$ 维的,而嵌入的 $1$ 维子集若稠密,就必须等于整个 $2$ 维流形)。

设 $\alpha:(-\pi/2,\pi/2)\to\mathbb{R}^2$ 由 $\alpha(t)=(t,\tan t)$ 定义。这是单射光滑浸入。其像集是函数 $f(t)=(\tan t)$ 的图,该函数定义在 $(-\pi/2,\pi/2)$ 上,故它是 $\mathbb{R}^2$ 的嵌入子流形。若我们把定义域改为 $(-\pi/2,\pi/2)$,但改用映射 $\widetilde{\alpha}:(-\pi/2,3\pi/2)\to\mathbb{R}^2$ 定义为
\begin{align*}
\widetilde{\alpha}(t)=
\begin{cases}
(t,\tan t), & -\pi/2<t<\pi/2,\\
(t,0), & \pi/2\leqslant t<3\pi/2.
\end{cases}
\end{align*}
这仍是单射光滑浸入(\cref{figure:ch05-immersed-manifold})。然而,其像 $S$ 在由 $\widetilde{\alpha}$ 诱导的拓扑下并非 $\mathbb{R}^2$ 的嵌入子流形,尽管它作为子集恰好就是嵌入子流形 $f$ 的图。原因在于 $\widetilde{\alpha}$ 是从开区间到 $S$ 上的同胚,而 $S$ 在子空间拓扑下不是流形,因为在点 $(\pi,0)$ 的任一邻域都包含来自图像的两部分。
\end{example}

\begin{figure}[ht]
\centering
\includegraphics[width=0.8\textwidth]{fig-05-8.png}
\caption{浸入子流形局部是嵌入的}\label{figure:ch05-immersed-manifold}
\end{figure}

注意\cref{example:ch05-immersed-examples}(c) 表明:浸入子流形 $S\subseteq M$ 不是流形这一断言,只有在给定拓扑与光滑结构的语境下才有意义。事实上,同一子集 $S$ 既可以是嵌入子流形又可以是浸入子流形,取决于我们对它赋予何种拓扑与光滑结构。(然而,后面的\cref{theorem:ch05-uniqueness} 将表明,如果 $S$ 是嵌入子流形,则它只容许一种使其成为嵌入子流形的光滑结构;而作为浸入子流形,它可容许许多不同的光滑结构。)

然而,一个子集不可能既是嵌入子流形,又在其上另加一个与子空间拓扑不同的拓扑而仍成为浸入子流形——除了最简单的情形外。下一个命题描述了唯一可能发生歧义的情形。

\begin{proposition}\label{proposition:ch05-immersed-vs-open}
设 $M$ 是光滑流形, $S\subseteq M$ 是浸入子流形,并设 $S$ 上的拓扑恰好就是子空间拓扑。则 $S$ 是嵌入子流形。
\end{proposition}

\begin{proof}
这直接由嵌入子流形与浸入子流形的定义推出:若包含映射 $S\hookrightarrow M$ 是单射浸入且同时是拓扑嵌入,则它是光滑嵌入。
\end{proof}

浸入子流形与嵌入子流形之间的差别反映在切片坐标卡的存在性上。回忆嵌入子流形容许切片坐标卡,使该子流形在每一这样的坐标卡中恰好是单个切片。对浸入子流形,这只在局部成立,且每次限制在浸入定义域的一个小邻域内时成立。

\begin{proposition}\label{proposition:ch05-immersed-slice-local}
设 $M$ 是光滑流形。若 $S\subseteq M$ 是浸入子流形,则 $S$ 的每一点都有 $M$ 中的邻域 $U$,使 $S\cap U$ 是 $U$ 的单个切片。
\end{proposition}

\begin{proof}
设 $S$ 是 $k$ 维浸入子流形, $p\in S$ 任意。由于包含映射 $\iota:S\hookrightarrow M$ 是浸入, \cref{theorem:ch04-local-embedding} 保证存在 $p$ 在 $S$ 中的邻域 $U_0$,使 $\iota|_{U_0}:U_0\to M$ 是嵌入。由于 $U_0$ 是 $S$ 的嵌入子流形, \cref{theorem:ch05-slice-criterion} 表明存在 $M$ 的坐标卡 $(V,\varphi)$ 使 $U_0$ 是 $V$ 中的 $k$-切片。然而, $S\cap V$ 未必等于 $U_0$,因为 $S$ 的其他部分可能进入 $V$。若发生这种情况, $\iota^{-1}(V)$ 在 $S$ 中将是含 $U_0$ 的一个连通分支的集合,它可能是不连通的,其他分支也可能映进 $V$。但若 $V$ 足够小,使得 $\iota^{-1}(V)$ 的含 $U_0$ 的分支就是 $U_0$ 本身,则 $S\cap V=U_0$。这可以做到,例如取 $V$ 充分小,使 $\iota^{-1}(V)$ 的其他分支在 $S$ 中与 $\overline{U_0}$ 不交。于是 $S\cap V=U_0$ 是 $V$ 的切片。
\end{proof}

一个自然的问题是:若 $S\subseteq M$ 是浸入子流形,是否存在使 $S$ 成为嵌入子流形的条件?下面两个命题给出部分答案。

\begin{proposition}\label{proposition:ch05-immersed-embedded-criteria}
设 $M$ 是光滑流形, $S\subseteq M$ 是浸入子流形。若下列任一条件成立,则 $S$ 实际上是 $M$ 的嵌入子流形。
\begin{enumerate}[(a)]
\item $S$ 是 $M$ 的开子集。
\item $S$ 是 $M$ 的闭子集。
\item $S$ 是紧的。
\end{enumerate}
\end{proposition}

\begin{proof}
若 $S$ 开,则由\cref{proposition:ch05-open-submanifolds} 知它是嵌入子流形。若 $S$ 闭或紧,则包含映射 $S\hookrightarrow M$ 是固有映射:若 $S$ 闭,这是显然的;若 $S$ 紧,则由\cref{proposition:appA-proper-criterion} 知它固有。无论哪种情形,由\cref{proposition:ch04-injective-immersion-embed} 知包含映射是嵌入。
\end{proof}

\end{document}